\subsection{Wie ist das Projekt in das eigene Unternehmen eingebettet?}
\label{subsec:Projekteinbettung}

\par Der \gls{adac} ist mit rund 23 Millionen Mitgliedern der führende Verkehrsclub Europas. Die Tochtergesellschaft \gls{ais} fungiert dabei als zentraler IT-Dienstleister und verantwortet die gesamte Infrastruktur sowie die digitale Transformation des Konzernverbunds.\endnote{Der Begriff Konzern ist hier nicht im juristischen Sinne (§ 18 AktG, § 290 HGB) zu verstehen, sondern bezeichnet das Gefüge bzw. Konglomerat aus Vereinen und Unternehmen unter dem Dach des \gls{ev}.} In einem Umfeld, das jährlich eine dreistellige Anzahl an IT-Projekten bewältigt, stellt das Programm \gls{r2h} das strategische Kernvorhaben dar. Ziel ist die Ablösung des 25 Jahre alten \gls{r3}-Systems durch \gls{s4}, um durch Prozessoptimierung und Automatisierung die langfristige Wettbewerbsfähigkeit sowie die regulatorische Konformität des Gesamtkonzerns sicherzustellen. Die Einbettung des Projekts in die \gls{ais} unterstreicht die Relevanz einer modernen IT-Plattform als Enabler für die Geschäftsmodelle des \gls{adac}.


\begin{comment}
\par Der \ac{adac} ist mit rund 23 Millionen Mitgliedern der führende Verkehrsclub Europas. Die Tochtergesellschaft \ac{ais} fungiert dabei als zentraler IT-Dienstleister und verantwortet die gesamte Infrastruktur sowie die digitale Transformation des Konzernverbunds. In einem Umfeld, das jährlich eine dreistellige Anzahl an IT-Projekten bewältigt, stellt das Programm \acl{r2h} (\ac{r2h}) das strategische Kernvorhaben dar. Ziel ist die Ablösung des 25 Jahre alten \ac{r3}-Systems durch \ac{s4}, um durch Prozessoptimierung und Automatisierung die langfristige Wettbewerbsfähigkeit sowie die regulatorische Konformität des Gesamtkonzerns sicherzustellen. Die Einbettung des Projekts in die \ac{ais} unterstreicht die Relevanz einer modernen IT-Plattform als Enabler für die Geschäftsmodelle des \ac{adac}.


\par Der \acl{adac} (\ac{adac}) ist mit rund 23 Millionen Mitgliedern der größte Verkehrsclub Europas mit Sitz in München. Die \acl{ais} (\acs{ais}) ist eine Tochtergesellschaft der \acl{se} (\acs{se}), eine der drei Säulen des \ac{adac}. Als IT-Dienstleister betreibt die \ac{ais} die gesamte IT-Infrastruktur des Konzert-Verbunds.\footnote{Der Begriff Konzern ist hier nicht im juristischen Sinne (§ 18 AktG, § 290 HGB) zu verstehen, sondern bezeichnet das Gefüge/Konglomerat aus Vereinen und Unternehmen unter dem Dach des ADAC.} Dies umfasst Entwicklung, Implementierung, Betrieb, Wartung und Decommissionierung sämtlicher Hard- und Softwarelösungen für die Geschäftsprozesse des \ac{adac} und seiner Tochtergesellschaften – teils in Eigenleistung, teils über Subunternehmen.

\par Pro Jahr führt die \acs{ais} eine dreistellige Anzahl von IT-Projekten durch, um die komplexe und vielseitige IT-Landschaft des \ac{adac} weiterzuentwickeln und an Marktanforderungen anzupassen. Die Projekte reichen von der Wartung bestehender Systeme über die Einführung neuer Software bis hin zur grundlegenden Modernisierung mit Implementierung neuer Technologien, etwa im Bereich Cloud-Computing und Künstliche Intelligenz. Ebenso vielfältig sind die passgenau gewählten Projektmethoden (klassisch, agil oder hybrid) und die Projektmanagement-Tools, welche die Planung, Überwachung und Steuerung unterstützen.

\par Ein zentrales Programm der digitalen Transformation ist \acl{r2h} (\ac{r2h}), das das bisherige \acl{erp}-System (\ac{erp}) \acl{r3} (\ac{r3}) schrittweise durch \acl{s4} (\ac{s4}) ersetzt.\footnote{Enterprise Resource Planning (\acs{erp}) umfasst die Planung, Steuerung und Verwaltung von Ressourcen und Prozessen zur Sicherstellung effizienter Wertschöpfung und optimierter Abläufe.} Die Transformationsreihe dient der Effizienzsteigerung durch Automatisierung und Prozessoptimierung. Ziel ist es, die Wettbewerbsfähigkeit des \ac{adac} langfristig durch moderne IT-Plattformen zu sichern und bestmögliche Unterstützung für sich stetig wandelnde technische, fachliche und regulatorische Anforderungen bereitzustellen.
\end{comment}